% answers_short.tex —— 1985 年·广东卷·文史类 · 简短答案速查
% 由 tools/make_short_answers.py 自 answers.tex 的 \gksolhead/\gkanswer 提取，
% 并与 1985/1985_数学真题_广东文史类_答案要点.md 逐题比对后生成，勿手改（改源 MD 后重跑）。

\begin{gknotice}
\textbf{使用说明}：本节只给最终答案便于快速核对，完整推导见解析版。答案为编写者独立推导，\textbf{非官方标准答案}。
\end{gknotice}

\gkpart{卷一填空}
\begin{itemize}[leftmargin=2.6em]
  \item[\textbf{1.}] $2$
  \item[\textbf{2.}] $64$
  \item[\textbf{3.}] $2\bigl(\cos\frac{11\pi}{6}+i\sin\frac{11\pi}{6}\bigr)$
  \item[\textbf{4.}] $\dfrac{x}{x^2-1}$（题设仅确定 $x\ne0,\pm1$ 时的函数值）
  \item[\textbf{5.}] $\dfrac19$
  \item[\textbf{6.}] $0$
  \item[\textbf{7.}] $\dfrac32$
  \item[\textbf{8.}] $8$（按棱的垂直方向计数；若“垂直”限于相交则为 4，见 U08）
  \item[\textbf{9.}] $\dfrac\pi3$ 弧度
  \item[\textbf{10.}] $\dfrac\pi2$
\end{itemize}

\gkpart{卷一选择}
\begin{itemize}[leftmargin=2.6em]
  \item[\textbf{1.}] E
  \item[\textbf{2.}] B
  \item[\textbf{3.}] E
  \item[\textbf{4.}] C
  \item[\textbf{5.}] D
  \item[\textbf{6.}] D
  \item[\textbf{7.}] C
  \item[\textbf{8.}] A
  \item[\textbf{9.}] B
  \item[\textbf{10.}] B
  \item[\textbf{11.}] U/PENDING；指数字形未能确认
  \item[\textbf{12.}] D
  \item[\textbf{13.}] B
  \item[\textbf{14.}] A
  \item[\textbf{15.}] D
\end{itemize}

\gkpart{卷二解答}
\begin{itemize}[leftmargin=2.6em]
  \item[\textbf{1.}] 见解析版
  \item[\textbf{2.}] $\displaystyle\sum_{j=1}^{n}(2j-1)^2=\frac{n(4n^2-1)}3$，$n\ge1$。
  \item[\textbf{3.}] (1) 焦点 $(2,0)$ 与短轴端点 $(0,-2)$ 均在 $l$ 上；(2) $\dfrac{8\sqrt2}{3}$；(3) 抛物线，方程为 $x^2+y^2+2xy+4x-4y-4=0$
  \item[\textbf{4.}] (1) $CE\perp AB_1$；(2) $\dfrac{2\sqrt3}{3}\,\mathrm{cm}$；(3) $\cos\alpha=\dfrac{\sqrt{10}}5$
\end{itemize}
